\documentclass[11pt]{article}
\usepackage[T1]{fontenc}
\usepackage[utf8]{inputenc}
\usepackage{lmodern}
\usepackage[a4paper,margin=2.35cm]{geometry}
\usepackage{amsmath,amssymb,amsthm,mathtools}
\usepackage{microtype,booktabs,array,longtable,multicol}
\usepackage[dvipsnames]{xcolor}
\usepackage{enumitem,listings,fancyhdr,titlesec,needspace}
\usepackage[colorlinks=true,linkcolor=MidnightBlue,citecolor=MidnightBlue,
 urlcolor=MidnightBlue]{hyperref}
\definecolor{Ink}{HTML}{16324F}
\definecolor{Accent}{HTML}{4E5C3D}
\definecolor{Wash}{HTML}{F1F4F6}
\titleformat{\section}{\Large\bfseries\color{Ink}}{\thesection}{0.7em}{}
\titleformat{\subsection}{\large\bfseries\color{Accent}}{\thesubsection}{0.7em}{}
\pagestyle{fancy}\fancyhf{}
\fancyhead[L]{\small\color{Ink}Prescribed coefficient-lattice minima}
\fancyhead[R]{\small Research draft}
\fancyfoot[C]{\thepage}
\setlength{\headheight}{14pt}
\setlength{\parskip}{0.32em}
\setlength{\parindent}{1.1em}
\emergencystretch=2em
\numberwithin{equation}{section}
\newtheorem{theorem}{Theorem}[section]
\newtheorem{lemma}[theorem]{Lemma}
\newtheorem{proposition}[theorem]{Proposition}
\newtheorem{corollary}[theorem]{Corollary}
\theoremstyle{definition}
\newtheorem{definition}[theorem]{Definition}
\newtheorem{example}[theorem]{Example}
\theoremstyle{remark}\newtheorem{remark}[theorem]{Remark}
\newcommand{\Z}{\mathbb Z}\newcommand{\N}{\mathbb N}
\newcommand{\Q}{\mathbb Q}\newcommand{\R}{\mathbb R}
\newcommand{\F}{\mathbb F}
\newcommand{\Latt}{\mathcal L}
\newcommand{\dgr}{\operatorname{d}}
\newcommand{\rank}{\operatorname{rank}}
\newcommand{\Span}{\operatorname{span}}
\newcommand{\supp}{\operatorname{supp}}
\newcommand{\diam}{\operatorname{diam}}
\newcommand{\norm}[1]{\left\lVert#1\right\rVert_1}
\newcommand{\pos}[1]{\left(#1\right)_+}
\lstset{basicstyle=\ttfamily\footnotesize,breaklines=true,columns=fullflexible,
 keepspaces=true,frame=single,rulecolor=\color{Ink!30},backgroundcolor=\color{Wash},
 showstringspaces=false,tabsize=4}
\hypersetup{pdftitle={Prescribing every successive L1-minimum of a coefficient lattice},
 pdfauthor={AI-assisted research draft prepared for Vladimir Reshetnikov},
 pdfsubject={Mixed-radix realization of O'Bryant's successive-minima conjecture}}
\title{\color{Ink}\bfseries Prescribing every successive\texorpdfstring{ $L^1$}{ L1}-minimum\par
 of a coefficient lattice\\[0.5em]
 \large Mixed-radix realization, rigid minimizing bases,\par
 and exact iterated-sumset growth}
\author{Research draft prepared for Vladimir Reshetnikov}
\date{20 September 2026}

\begin{document}
\maketitle
\begin{abstract}
Let $\Latt(A)$ be the integer vectors perpendicular to both the all-ones
vector and a vector listing the distinct integers of a finite set $A$.
We prove that every nondecreasing list of even integers at least four occurs
as the complete list of successive $L^1$-minima of such a lattice, with exactly
two more elements in $A$ than entries in the list.  For prescribed half-minima
$2\le h_1\le\cdots\le h_r$, an explicit witness is
\[
 A=\{0,1,h_1,h_1h_2,\ldots,h_1\cdots h_r\}.
\]
A stronger theorem permits the multipliers in any order.  Its proof gives an
integral basis of carry relations and bounds each basis coordinate of every
short vector.  It determines the entire subgroup generated by short relations
and all successive minimizing bases.  This resolves the substantive
realization conjecture following Lemma~1 of O'Bryant's
\emph{On Nathanson's Triangular Number Phenomenon}, while correcting its
impossible endpoint $h_1=1$.  We also obtain unique mixed-radix normal forms,
a complete-intersection presentation, an exact generating function for every
iterated sumset, and the sharp onset of linear growth.  All arguments are
self-contained; exact-arithmetic verification programs and their outputs
accompany the article.
\end{abstract}

\noindent\colorbox{Wash}{\parbox{\dimexpr\linewidth-2\fboxsep\relax}{\small
\textbf{Status.} This is an AI-assisted mathematical research draft, not an
independently refereed or proof-assistant-verified paper.  The all-parameter
claims are supported by the proofs below, not by extrapolation from tests.
The source still states the target as a conjecture in the version inspected;
no later solution was located in the searches recorded with this package.
That search does not establish priority.}}

\clearpage
\tableofcontents
\clearpage

\section{The problem and the precise claim}

The inverse problem studied here asks how freely one can prescribe the
sizes at which linearly independent additive relations first appear in a set
of integers.  A relation between two multisets of the same cardinality is
encoded by subtracting their multiplicity vectors.  Its $L^1$-norm is twice
the number of terms remaining on either side after cancellation.  Prescribing
successive minima therefore means prescribing successive independent
collision thresholds, not merely finding one short relation.

The explicit target is the unnumbered conjecture immediately after Lemma~1
in Section~1.1, page~2, of O'Bryant's paper~\cite{obryant}.  That lemma handles
the first two minima; the following conjecture asks for a full prescribed
list.  The displayed lower endpoint is $h_1=1$.  Section~\ref{sec:necessary}
proves that this endpoint cannot occur for a set of distinct integers.
The realization question on its admissible domain is:
\begin{quote}
Given $r\ge1$ and integers $2\le h_1\le\cdots\le h_r$, is there an
$(r+2)$-element set of integers whose coefficient lattice has successive
$L^1$-minima $2h_1,\ldots,2h_r$?
\end{quote}
The answer is affirmative.  We do not stop at the endpoint obstruction:
we construct a witness for \emph{every} admissible list.

\subsection{What is and is not being resolved}

The central result is a full realization theorem, not a statement about a
fixed number of minima, a generic parameter range, or a bounded search.
Repeated minima and arbitrarily large gaps between consecutive minima are
allowed.  The witness is normalized by $\min A=0$ and $1\in A$.
Its lattice has a genuine $\Z$-basis attaining all the minima.

The additional sumset formulas apply to the constructed multiplicative-chain
sets.  They do not classify all sets with the same minima, nor all attainable
sumset cardinalities.  The larger range-of-sizes problem is discussed, for
example, in Nathanson's work~\cite{nathanson}.  Section~\ref{sec:limits}
gives two sets with the same complete minimum list but different later
sumset sizes.  Thus the special-family formulas must not be promoted to
formulas for arbitrary coefficient lattices.

\subsection{Selection and exclusion audit}

Before selecting a problem, a single call to an operating-system-backed
random number generator selected entry~1 from a fixed catalogue of 96 areas:
\emph{additive combinatorics}.  The complete catalogue and the recorded result
are included in \texttt{data/area\_selection.json}; the catalogue is also
printed in Appendix~\ref{app:areas}.  There was no redraw.

All 71 entries in the supplied \texttt{manifest(1).tex} were inspected as
exclusions.  None describes this coefficient-lattice realization problem.
The comparison uses the manifest's descriptions; it is not a re-verification
of the proofs in those other packages.  The target and the distinction from
the excluded topics are recorded in \texttt{notes/exclusion\_audit.md}.

\section{Coefficient lattices and necessary conditions}
\label{sec:necessary}

Fix a finite integer set
\[
 A=\{a_0<a_1<\cdots<a_{k-1}\},\qquad k\ge2.
\]
Coordinates are numbered from zero throughout this article.

\begin{definition}
The \emph{coefficient lattice} of $A$ is
\begin{equation}\label{eq:lattice}
 \Latt(A)=\left\{c\in\Z^k:
     \sum_{i=0}^{k-1}c_i=0,\quad
     \sum_{i=0}^{k-1}a_i c_i=0\right\}.
\end{equation}
For $c\in\Latt(A)$ put $c_i^+=\max(c_i,0)$,
$c_i^-=\max(-c_i,0)$, and
\begin{equation}\label{eq:degree}
 \dgr(c)=\sum_i c_i^+=\sum_i c_i^-=\frac12\norm{c}.
\end{equation}
For $1\le j\le k-2$, its $j$th successive minimum is
\begin{equation}\label{eq:minima}
 \lambda_j(A)=\min\left\{T\ge0:
 \dim_\Q\Span_\Q\{c\in\Latt(A):\norm{c}\le T\}\ge j\right\}.
\end{equation}
\end{definition}
The two rows defining~\eqref{eq:lattice} are independent because the $a_i$
are not all equal.  Their rational kernel has dimension $k-2$.  A rational
basis can be multiplied by denominators to give independent integer
vectors, so $\Latt(A)$ has rank $k-2$.  Hence all the minima in
\eqref{eq:minima} exist.  Using $\R$-linear instead of $\Q$-linear
independence gives the same definition for these rational vectors.

\begin{proposition}[The complete elementary obstruction]
\label{prop:necessary}
Every successive minimum of $\Latt(A)$ is an even integer at least four.
\end{proposition}
\begin{proof}
Equation~\eqref{eq:degree} makes every lattice norm an even integer.  If a
nonzero vector had degree one, its positive and negative parts would each
be a coordinate unit vector.  Thus it would equal $e_i-e_j$ for $i\ne j$.
The second relation in~\eqref{eq:lattice} would imply $a_i=a_j$, contrary
to distinctness.  A nonzero relation therefore has degree at least two.
\end{proof}

Consequently, the conjecture with the literal endpoint $h_1=1$ is false.
Replacing that endpoint by two is necessary, and the next section proves
that it is sufficient.  The adjacent two-minimum lemma in the source already
uses the endpoint two~\cite{obryant}; the endpoint discrepancy is stated
explicitly here rather than silently repaired.

Translation and nonzero scaling do not change the lattice: if
$b_i=u+v a_i$ with $v\ne0$, then, on $\sum c_i=0$,
$\sum b_i c_i=v\sum a_i c_i$.  A negative scaling merely reverses the
increasing coordinate order.  The construction below therefore also gives
strictly positive witnesses after adding one to every element.

\section{A full realization theorem}
\label{sec:main}

It is useful to prove more than the ordered-parameter statement.  Let
$q_1,\ldots,q_r\ge2$ be integers in \emph{any} order.  Define
\begin{equation}\label{eq:construction}
 a_0=0,\qquad a_1=1,\qquad a_{j+1}=q_j a_j\quad(1\le j\le r),
 \qquad A(q)=\{a_0,\ldots,a_{r+1}\}.
\end{equation}
The elements are strictly increasing.  Write
\begin{equation}\label{eq:carry}
 v_j=(q_j-1)e_0-q_j e_j+e_{j+1}\qquad(1\le j\le r).
\end{equation}
This is the carry relation
\[
 \underbrace{a_j+\cdots+a_j}_{q_j\text{ terms}}
 =a_{j+1}+\underbrace{0+\cdots+0}_{q_j-1\text{ terms}}.
\]

\Needspace{17\baselineskip}
\begin{theorem}[Mixed-radix spectrum theorem]
\label{thm:spectrum}
For $A(q)$ in~\eqref{eq:construction}, the following statements hold.
\begin{enumerate}[label=\textup{(\roman*)},itemsep=0.25em]
\item The vectors $v_1,\ldots,v_r$ are a $\Z$-basis of $\Latt(A(q))$.
\item If $c=\sum_{j=1}^r z_jv_j$, then, for every $j$,
\begin{equation}\label{eq:coordinatebound}
 q_j|z_j|\le\dgr(c).
\end{equation}
\item For every real $H\ge0$, the subgroup generated by relations of degree
at most $H$ is exactly
\begin{equation}\label{eq:filtration}
 G_H:=\left\langle c\in\Latt(A(q)):\dgr(c)\le H\right\rangle_\Z
       =\bigoplus_{q_j\le H}\Z v_j.
\end{equation}
\item If $q_{(1)}\le\cdots\le q_{(r)}$ is the sorted list of the radices,
then
\begin{equation}\label{eq:spectrum}
 (\lambda_1(A(q)),\ldots,\lambda_r(A(q)))
       =(2q_{(1)},\ldots,2q_{(r)}).
\end{equation}
\end{enumerate}
\end{theorem}

The upper bounds are supplied by the explicit carries.  The substantive
point is to rule out short integer combinations that could create extra
independent directions earlier.  The coordinate estimate does precisely that.

\subsection{The integral basis}

\begin{lemma}\label{lem:basis}
The carries form a $\Z$-basis, and their coordinates are given by
\begin{equation}\label{eq:coordinateformula}
 z_j=-\frac{1}{a_{j+1}}\sum_{i=1}^{j}a_i c_i.
\end{equation}
\end{lemma}
\begin{proof}
Each $v_j$ has coordinate sum zero and weighted sum
$-q_j a_j+a_{j+1}=0$, so lies in the lattice.  Its highest nonzero
coordinate is $j+1$, with coefficient one.  Reading from the highest
coordinate down shows that the $v_j$ are independent.

For an arbitrary $c\in\Latt(A(q))$, subtract $c_{r+1}v_r$.  This kills
coordinate $r+1$ without changing the lattice constraints.  Continue
backwards, killing coordinate $j+1$ with an integer multiple of $v_j$.
The remainder is supported on coordinates zero and one.  Since $a_0=0$
and $a_1=1$, its weighted-sum constraint kills coordinate one, and its
coordinate-sum constraint then kills coordinate zero.  This proves integral,
not just rational, spanning.

For the coordinate formula, calculate the partial weighted sum of each
basis vector through coordinate $j$.  The result is zero for $v_\ell$ with
$\ell<j$, zero for $\ell>j$, and $-a_{j+1}$ for $v_j$.
Linearity gives~\eqref{eq:coordinateformula}.
\end{proof}

\subsection{The inequality excluding hidden short relations}

\begin{lemma}[A cut-coordinate bound]\label{lem:bound}
For every $c\in\Latt(A(q))$ and every $1\le j\le r$,
$q_j|z_j|\le\dgr(c)$.
\end{lemma}
\begin{proof}
Set
\[
 P_j=\sum_{i=1}^j a_i c_i^+,
 \qquad N_j=\sum_{i=1}^j a_i c_i^-.
\]
All weights in these sums lie between zero and $a_j$, while each positive
or negative mass is at most $\dgr(c)$.  Thus
$0\le P_j,N_j\le a_j\dgr(c)$, and consequently
\[
 a_{j+1}|z_j|
 =\left|\sum_{i=1}^j a_i c_i\right|
 =|P_j-N_j|
 \le a_j\dgr(c).
\]
Division by $a_j>0$ and $a_{j+1}=q_j a_j$ proves the claim.
\end{proof}

\begin{proof}[Completion of Theorem~\ref{thm:spectrum}]
If $\dgr(c)\le H$ and $q_j>H$, the bound forces $z_j=0$.
Hence every generator of $G_H$ lies in the right side of
\eqref{eq:filtration}.  Conversely, $\dgr(v_j)=q_j$, so every basis
vector on that right side is itself one of the allowed generators.
This proves equality of subgroups.  Their rational rank is
\begin{equation}\label{eq:rankprofile}
 \dim_\Q\Span_\Q\{c:\dgr(c)\le H\}=\#\{j:q_j\le H\}.
\end{equation}
The definition of successive minima now gives~\eqref{eq:spectrum}.
\end{proof}

\begin{corollary}[Exact admissibility classification]\label{cor:realization}
For $r\ge1$, a nondecreasing list $(\ell_1,\ldots,\ell_r)$ is the full
successive-$L^1$-minimum list of the coefficient lattice of a set of
$r+2$ distinct integers if and only if every $\ell_j$ is even and
$\ell_1\ge4$.  For $\ell_j=2h_j$, one witness is
\[
 \boxed{\ A=\{0,1,h_1,h_1h_2,\ldots,h_1\cdots h_r\}.\ }
\]
\end{corollary}
\begin{proof}
Necessity is Proposition~\ref{prop:necessary}; sufficiency is
Theorem~\ref{thm:spectrum} with $q_j=h_j$.
\end{proof}
For $r=0$, the empty minimum list is realized by $\{0,1\}$, whose
coefficient lattice is zero.  This is the only rank-zero convention used.

\section{Rigidity, short-vector enumeration, and a second lower-bound argument}

\subsection{Only the carry vectors introduce a direction at its first level}

Let
\[
 G_{<Q}=\left\langle c\in\Latt(A(q)):\dgr(c)<Q\right\rangle_\Z
       =\bigoplus_{q_j<Q}\Z v_j.
\]
The filtration theorem is stronger than a list of minima.  In fact, the
vectors that first leave one of its layers are completely rigid.

\begin{proposition}[Entrance rigidity]\label{prop:rigidity}
For any radix value $Q$ that occurs in the list,
\begin{equation}\label{eq:rigidity}
 \{c\in\Latt(A(q)):\dgr(c)\le Q,\ c\notin G_{<Q}\}
      =\{v_j,-v_j:q_j=Q\}.
\end{equation}
\end{proposition}
\begin{proof}
Take a vector on the left.  Some basis coordinate $z_j$ is nonzero with
$q_j\ge Q$.  The inequalities
$q_j|z_j|\le\dgr(c)\le Q$ imply $q_j=Q$, $|z_j|=1$ and
$\dgr(c)=Q$.  Change the sign so that $z_j=1$.
In the notation of Lemma~\ref{lem:bound},
$N_j-P_j=Qa_j$, while $N_j\le Qa_j$ and $P_j\ge0$.
Therefore $N_j=Qa_j$ and $P_j=0$.

Equality in $N_j\le a_j\sum_{i=1}^j c_i^-\le a_j\dgr(c)$,
with strictly increasing positive weights, forces all negative mass to be
at coordinate $j$: $c_j=-Q$.  Every positive nonzero-weight coordinate must
then be beyond $j$.  Their total weighted sum is $Qa_j=a_{j+1}$.
The smallest available weight is $a_{j+1}$, so exactly one unit of positive
mass occurs there, and none farther to the right.  The other $Q-1$ units
of positive mass are at coordinate zero.  Thus $c=v_j$.
The reverse inclusion follows directly from the basis and its degrees.
\end{proof}

\begin{corollary}[All minimizing bases]\label{cor:bases}
Let $m_Q=\#\{j:q_j=Q\}$.  Every ordered independent list
$c_1,\ldots,c_r$ satisfying $\norm{c_i}=\lambda_i(A(q))$ consists of the
carry vectors, with arbitrary signs and arbitrary permutations among equal
radices.  Every such list is a $\Z$-basis.  The number of these ordered lists is
\begin{equation}\label{eq:basiscount}
 2^r\prod_Q m_Q!.
\end{equation}
\end{corollary}
\begin{proof}
At the smallest degree, Proposition~\ref{prop:rigidity} gives exactly the
signed carries at that degree.  Independence forces the whole block to use
each direction once.  Induct over the distinct degree values.  Previous
blocks span $G_{<Q}$; an independent vector of degree $Q$ must leave that
space, so the proposition applies again.  Lemma~\ref{lem:basis} and elementary
sign and permutation counting finish the proof.
\end{proof}

\subsection{An exact, bounded search in carry coordinates}

For $r\ge1$, the norm can be written explicitly as
\begin{align}\label{eq:normcoordinates}
 2\dgr\!\left(\sum_jz_jv_j\right)
  ={}&\left|\sum_{j=1}^r(q_j-1)z_j\right|+q_1|z_1|\notag\\
     &+\sum_{j=2}^r|z_{j-1}-q_jz_j|+|z_r|.
\end{align}
A degree-$H$ ball can therefore be enumerated by trying precisely the
integer box
\begin{equation}\label{eq:box}
 |z_j|\le\left\lfloor\frac{H}{q_j}\right\rfloor
 \quad(1\le j\le r)
\end{equation}
and retaining the vectors that pass~\eqref{eq:normcoordinates}.
At most
$\prod_j(2\lfloor H/q_j\rfloor+1)$ candidates need to be tried.
This is an exact finite algorithm, not a claim of polynomial complexity in
$r$.  The box is an \emph{outer bound}, not a description of the norm ball.
Similarly, $G_H$ in~\eqref{eq:filtration} is a subgroup generated by a ball;
it is not the ball itself and contains vectors of arbitrarily large norm
as soon as it is nonzero.

\subsection{A short alternative proof when the radices are ordered}

The original realization claim can also be checked without the coordinate
formula.  Suppose $q_1\le\cdots\le q_r$, take $0\ne c\in\Latt(A(q))$,
and let $m$ be its highest nonzero coordinate.  Necessarily $m\ge2$.
Change the sign so that $c_m<0$.  The weighted negative part is at least
$a_m$, while the weighted positive part is at most
$\dgr(c)a_{m-1}$.  Since those weighted sums are equal,
\begin{equation}\label{eq:highest}
 \dgr(c)\ge a_m/a_{m-1}=q_{m-1}.
\end{equation}
Thus relations of degree less than $q_j$ are supported on coordinates
$0,\ldots,j$.  The relation space on these $j+1$ coordinates has dimension
$j-1$, so it cannot contain $j$ independent vectors.  This gives
$\lambda_j\ge2q_j$; the first $j$ carries give the reverse inequality.
Unlike the cut-coordinate proof, this particular support argument uses
ordered radices.

\section{Mixed-radix normal forms and all iterated sumsets}
\label{sec:normal}

Put
\begin{equation}\label{eq:PD}
 P=a_{r+1}=\prod_{j=1}^r q_j,\qquad
 D=\sum_{j=1}^r(q_j-1),\qquad
 \mathcal D=\prod_{j=1}^r\{0,1,\ldots,q_j-1\}.
\end{equation}
For a digit tuple $d=(d_1,\ldots,d_r)\in\mathcal D$, write
$s(d)=\sum_jd_j$.  Let $hA$ denote sums of exactly $h$ elements of $A$,
with repetition allowed, and let $0A=\{0\}$.

\begin{lemma}[Unique digits and minimum number of summands]
\label{lem:normal}
Every integer $n\ge0$ has a unique expansion
\begin{equation}\label{eq:digits}
 n=tP+\sum_{j=1}^r d_j a_j,
 \qquad t\ge0,\quad d\in\mathcal D.
\end{equation}
The minimum number of nonzero elements of $A(q)$ needed to represent $n$
is $t+s(d)$.  Consequently,
\begin{equation}\label{eq:membership}
 n\in hA(q)\quad\Longleftrightarrow\quad t+s(d)\le h.
\end{equation}
\end{lemma}
\begin{proof}
Repeated Euclidean division by $q_1,q_2,\ldots,q_r$ gives the digits and
then the quotient $t$.  Conversely, reduction modulo $q_1$ determines $d_1$;
subtraction and division determine $d_2$, and so on.  This proves uniqueness.

Start with any representation of $n$ by nonzero elements of $A(q)$.
Beginning at $a_1$, replace every $q_j$ copies of $a_j$ by one $a_{j+1}$.
Process indices in increasing order.  Later replacements never reintroduce
a smaller denomination, and each replacement reduces the number of nonzero
summands by $q_j-1$.  The final representation has bounded digits and so is
the unique expansion~\eqref{eq:digits}.  It follows that its $t+s(d)$
summands are minimal.  Zero summands may then be added to obtain exactly
$h$ terms, proving~\eqref{eq:membership}.
\end{proof}

No general assertion that a greedy coin algorithm is optimal is being
used: optimality was proved for this divisibility chain by explicit carries.
In particular, there is no floating-point comparison and no search in the
membership test.

\begin{theorem}[Exact sumset enumeration]\label{thm:sumsets}
For every integer $h\ge0$,
\begin{equation}\label{eq:digitscount}
 |hA(q)|=\sum_{d\in\mathcal D}\max\{0,h-s(d)+1\}.
\end{equation}
The ordinary generating function is
\begin{equation}\label{eq:gf}
 \boxed{\ \sum_{h\ge0}|hA(q)|z^h
 =\frac{\displaystyle\prod_{j=1}^r(1+z+\cdots+z^{q_j-1})}{(1-z)^2}
 =\frac{\displaystyle\prod_{j=1}^r(1-z^{q_j})}{(1-z)^{r+2}}.\ }
\end{equation}
More precisely, the bivariate enumerator is
\begin{equation}\label{eq:bivariate}
 \sum_{h\ge0}\sum_{n\in hA(q)}z^h w^n
 =\frac{\displaystyle\prod_{j=1}^r
       \left(\sum_{d=0}^{q_j-1}(zw^{a_j})^d\right)}
       {(1-z)(1-zw^P)}.
\end{equation}
\end{theorem}
\begin{proof}
For each digit tuple, exactly the quotients
$0\le t\le h-s(d)$ are allowed by~\eqref{eq:membership}.
Distinct digit/quotient pairs give distinct integers, proving
\eqref{eq:digitscount}.

Equivalently, every pair $(h,n)$ on the left of~\eqref{eq:bivariate}
is represented uniquely by $d\in\mathcal D$, $t\ge0$, and the number
$b=h-t-s(d)\ge0$ of zero summands.  Their weights are
$z^b$, $(zw^P)^t$, and $\prod_j(zw^{a_j})^{d_j}$.
Summing these independent choices gives~\eqref{eq:bivariate}.
Setting $w=1$ and summing the finite geometric series gives~\eqref{eq:gf}.
All operations are valid as formal power series in $z$.
\end{proof}

\section{Consequences for growth and reconstruction}

\subsection{A finite inclusion--exclusion formula}

For an integer $m$ and nonnegative integer $b$, use
$\binom{m}{b}=0$ when $m<b$.  This is a coefficient-counting convention,
not the polynomial continuation of binomial coefficients to negative $m$.
Expanding the numerator of~\eqref{eq:gf} gives
\begin{equation}\label{eq:inclusion}
 |hA(q)|=\sum_{S\subseteq\{1,\ldots,r\}}(-1)^{|S|}
 \binom{h-\sum_{j\in S}q_j+r+1}{r+1}.
\end{equation}
This uses $2^r$ terms regardless of the magnitude of $P$.
It is therefore complementary to direct digit enumeration, which has $P$
digit tuples, and to polynomial multiplication, whose output degree is $D$.

When the sorted radices satisfy $q_{(1)}<q_{(2)}$, the interval
$q_{(1)}\le h<q_{(2)}$ has exactly one active nonempty singleton in
\eqref{eq:inclusion}, and hence
\[
 |hA(q)|=\binom{h+r+1}{r+1}
          -\binom{h-q_{(1)}+r+1}{r+1}.
\]
Below the smallest radix there are no collisions at all.  These initial
regimes agree with O'Bryant's general Theorem~2~\cite{obryant}; they are
not claimed as new general sumset theorems here.  If the smallest radix
$Q$ occurs $m$ times, then at $h=Q$ the deficit from
$\binom{Q+r+1}{r+1}$ is exactly $m$.

\subsection{Discrete curvature and the exact linearity threshold}

Define the digit-sum polynomial
\begin{equation}\label{eq:digitpoly}
 C_q(z)=\prod_{j=1}^r(1+z+\cdots+z^{q_j-1})
       =\sum_{s=0}^D p_s z^s.
\end{equation}
Here $p_s$ counts digit tuples of sum $s$.  Every integer $s$ between zero
and $D$ occurs: this follows inductively because the sum of two integer
intervals is an integer interval.  Thus $p_s>0$ throughout that interval.
The involution $d_j\mapsto q_j-1-d_j$ gives $p_s=p_{D-s}$.
Moreover,
\begin{equation}\label{eq:digitmoments}
 \sum_s p_s=P,\qquad \sum_s s p_s=\frac{PD}{2}.
\end{equation}
The second identity follows either by the same involution or by summing
one digit coordinate at a time.  In particular, $PD$ is even.

Writing $N(h)=|hA(q)|$ and $N(-1)=N(-2)=0$, equation~\eqref{eq:gf} yields
\begin{equation}\label{eq:curvature}
 N(h)-2N(h-1)+N(h-2)=p_h\qquad(h\ge0),
\end{equation}
where $p_h=0$ for $h>D$.  The first differences increase to $P$, and their
successive increases are the symmetric digit-sum distribution.

\begin{theorem}[Sharp onset of affine growth]\label{thm:affine}
If $r\ge1$, then for every $h\ge D-1$,
\begin{equation}\label{eq:affine}
 |hA(q)|=P(h+1)-\frac{PD}{2}.
\end{equation}
The least nonnegative index from which this formula holds permanently is
$D-1$.  If $D\ge2$, its failure immediately before that index is exactly
\begin{equation}\label{eq:preboundary}
 |(D-2)A(q)|=P(D-1)-\frac{PD}{2}+1.
\end{equation}
\end{theorem}
\begin{proof}
In~\eqref{eq:digitscount}, once $h\ge D-1$ every term $h-s(d)+1$ is
nonnegative; the possible boundary term at $s(d)=D$ is zero.  Remove the
maximum and use~\eqref{eq:digitmoments} to obtain~\eqref{eq:affine}.
At $h=D-2$, precisely one digit tuple has $s(d)=D$, and its untruncated
contribution would be $-1$.  All other untruncated contributions are
nonnegative.  Truncation therefore adds exactly one, giving
\eqref{eq:preboundary}.  If $D=1$, the formula already holds at $h=0$,
which is the least allowed index.
\end{proof}
For $r=0$ the separate convention is $N(h)=h+1$, valid for all $h\ge0$.
The denominator $(1-z)^2$ in the first fraction of~\eqref{eq:gf} is
reduced: its numerator takes the nonzero value $P$ at $z=1$.

\subsection{Permuting the radices and recovering their multiset}

Although the order of the radices changes the intermediate elements of
$A(q)$, it changes neither the successive minima nor any sumset size.
The first assertion is Theorem~\ref{thm:spectrum}; the second is the
symmetric product in~\eqref{eq:gf}.  The actual sumsets can change, as can
the coefficients in the carry basis.

\begin{proposition}[Reconstruction within the chain family]
\label{prop:reconstruct}
The complete sequence $(|hA(q)|)_{h\ge0}$ determines the multiset of
radices of a multiplicative-chain set.
\end{proposition}
\begin{proof}
The value at $h=1$ determines $r=|A(q)|-2$.  Multiply the generating series
by $(1-z)^{r+2}$ to recover $F(z)=\prod_j(1-z^{q_j})$.
If $F\ne1$, the least positive exponent with nonzero coefficient is the
smallest remaining radix $m$, and its coefficient is minus the multiplicity
of $m$.  Remove the corresponding factor $(1-z^m)$ to that multiplicity
and repeat.  The process terminates after the known number $r$ of factors.
\end{proof}
This is not an inverse theorem for arbitrary integer sets: membership in
the chain family is essential to the reconstruction argument.

\section{The associated binomial ideal}
\label{sec:algebra}

The sumset proof is elementary and independent of this section.  An
algebraic formulation nevertheless explains why its generating function
has such a simple product numerator.

Let $K$ be any field, and put $R=K[x_0,\ldots,x_{r+1}]$.  Give each
$x_i$ ordinary degree one.  Define
\[
 \varphi:R\longrightarrow K[u,t],\qquad x_i\longmapsto u t^{a_i}.
\]
The degree-$h$ image has one basis monomial $u^h t^n$ for each
$n\in hA(q)$, so its dimension is $|hA(q)|$.
For $1\le j\le r$ set
\begin{equation}\label{eq:binomials}
 f_j=x_j^{q_j}-x_0^{q_j-1}x_{j+1}.
\end{equation}

\begin{theorem}[Carry presentation]\label{thm:ideal}
The kernel of $\varphi$ is the ideal $(f_1,\ldots,f_r)$.
A $K$-basis of the quotient consists of the monomials
\begin{equation}\label{eq:monomials}
 x_0^b x_{r+1}^t\prod_{j=1}^r x_j^{d_j},
 \qquad b,t\ge0,\quad 0\le d_j<q_j.
\end{equation}
For lexicographic order
$x_1>x_2>\cdots>x_{r+1}>x_0$, the $f_j$ form a Gr\"obner basis with
leading monomials $x_j^{q_j}$.  They are also a regular sequence.
\end{theorem}
\begin{proof}
Every $f_j$ lies in the kernel because $q_j a_j=a_{j+1}$.
Replacing $x_j^{q_j}$ by $x_0^{q_j-1}x_{j+1}$, in increasing order of $j$,
reduces every monomial to one of~\eqref{eq:monomials}.  The replacement
preserves the ordinary degree and its weighted exponent.
Distinct normal monomials have distinct images: equality of their
$t$-exponents gives equality of their digits and top quotients by
Lemma~\ref{lem:normal}, and equality of their $u$-exponents then gives
equality of $b$.

A polynomial in $\ker\varphi$ reduces modulo the $f_j$ to a linear
combination of normal monomials.  Their distinct image monomials are
linearly independent in $K[u,t]$, so every coefficient in that remainder
is zero.  This proves both the kernel assertion and the quotient basis.

In the stated lexicographic order, each carry replacement reduces a
monomial: its first changed variable is $x_j$ and its exponent decreases.
Division by the $f_j$ therefore produces a remainder containing no
$x_j^{q_j}$.  The kernel argument says that an element of the ideal has
zero remainder, which is the Gr\"obner-basis property.  It also follows
from the fact that the leading monomials are pairwise relatively prime,
but that criterion is not needed here.

For regularity, impose the equations in the order $f_1,f_2,\ldots$.
After imposing $f_1,\ldots,f_m$, the quotient is a free module over
$K[x_0,x_{m+1},\ldots,x_{r+1}]$, with basis
$\prod_{j=1}^m x_j^{d_j}$, $0\le d_j<q_j$.
This follows inductively from division by the monic polynomial
$x_m^{q_m}-x_0^{q_m-1}x_{m+1}$.
The next polynomial $f_{m+1}$ is a nonzero element of that polynomial
base ring, an integral domain.  Multiplication by it is injective on the
free module.  Thus it is not a zero divisor after the preceding equations.
All the polynomials have positive degree, so the resulting ideal is proper.
This proves regularity.
\end{proof}

Thus the image ring is a complete intersection in the concrete sense of a
polynomial ring modulo the displayed regular sequence.  It is a free module
of rank $P$ over $K[x_0,x_{r+1}]$.  Counting the degrees of the normal
monomials recovers~\eqref{eq:gf} without invoking any general formula for
Hilbert series.  In particular, we did not infer ideal generation merely
from lattice-basis generation; that inference is not valid in general.

\section{Examples and limits of the construction}
\label{sec:limits}

\subsection{A five-element example}

For $q=(2,3,5)$, the construction is
\[
 A=\{0,1,2,6,30\},
\]
with carry basis
\[
 \begin{pmatrix}
 1&-2&1&0&0\\
 2&0&-3&1&0\\
 4&0&0&-5&1
 \end{pmatrix}.
\]
The successive norms are $4,6,10$.  At degree bounds $H<2$, $2\le H<3$,
$3\le H<5$, and $H\ge5$, the ranks generated by short relations are,
respectively, $0,1,2,3$.  Here
\[
 C_q(z)=1+3z+5z^2+6z^3+6z^4+5z^5+3z^6+z^7,
\]
and the sumset sizes begin
\[
 1,\ 5,\ 14,\ 29,\ 50,\ 76,\ 105,\ 135,\ 165,\ldots.
\]
Since $P=30$ and $D=7$, the exact affine formula is
$|hA|=30h-75$ for $h\ge6$; at $h=5$ it underestimates by exactly one.
The bivariate formula distinguishes the actual sums, while this ordinary
sequence records only their number.

\subsection{Repeated minima and large gaps}

For $q=(2,2,4,7)$,
\[
 A=\{0,1,2,4,16,112\},\qquad
 (\lambda_1,\lambda_2,\lambda_3,\lambda_4)=(4,4,8,14).
\]
Exactly two independent directions enter at degree two, but no direction
enters at degree three.  The number of ordered successive minimizing bases
is $2^4\cdot2!=32$.  Replacing $7$ by any larger integer postpones the
last independent relation to that degree without changing the earlier
minima.

For $q=(4,9)$, the witness is $\{0,1,4,36\}$ with successive norms $8,18$.
Inclusion--exclusion gives, for every $h\ge0$,
\[
 |hA|=\binom{h+3}{3}-\binom{h-1}{3}
      -\binom{h-6}{3}+\binom{h-10}{3}.
\]
The zero convention on binomial coefficients is essential for small $h$.
This example gives two independently chosen thresholds by the same
construction that works for an arbitrarily long list.

\subsection{Equal minima need not determine equal sumset sizes}

Consider
\[
 A=\{0,1,3,9\},\qquad B=\{0,1,5,7\}.
\]
Both have successive norms $(6,6)$.  This follows for $A$ from the theorem.
For $B$, its ten unordered two-term sums are
\[
 0,1,2,5,6,7,8,10,12,14,
\]
so there is no nonzero relation of degree at most two.  On the other hand,
\[
 (-2,2,1,-1),\qquad (0,1,-3,2)
\]
are independent coefficient relations of degree three.  This proves its
minimum list without assuming any unproved classification.

Nevertheless, $|4A|=27$ while $|4B|=26$.  The first value follows from
\eqref{eq:gf}.  To check the second directly, take $l$ summands from
$\{5,7\}$ and fill the remaining $4-l$ positions with zeros and ones.
For $l=0,1,2,3,4$ the resulting sets are, respectively,
\[
 [0,4],\quad [5,10],\quad [10,16],\quad [15,22],
 \quad\{20,22,24,26,28\}.
\]
Their union is $[0,22]\cup\{24,26,28\}$, of size 26.
Thus even the complete successive-minimum list does not specify a general
sumset sequence.

\subsection{The diameter is controlled, not optimized}

The explicit witness has diameter $P=\prod_jh_j$.  This can be much larger
than necessary.  If all $r$ half-minima are two, the interval
\[
 I=\{0,1,\ldots,r+1\}
\]
also works: its $r$ independent second-difference relations
$e_{j-1}-2e_j+e_{j+1}$ have norm four, the universal lower bound.
Every $(r+2)$-element integer set has diameter at least $r+1$, so this
interval has minimum possible diameter for that list.  The chain witness
instead has diameter $2^r$.

There is also an elementary general lower bound.  Suppose $k=r+2$ and the
first half-minimum is $h_1$.  At $s=h_1-1$, all unordered $s$-term sums are
distinct; otherwise their difference would be a relation of degree at most
$s$.  A set of diameter $M$ has $sA$ in an integer interval of length $sM$.
Therefore
\begin{equation}\label{eq:diameter}
 M\ge\max\left\{k-1,
 \left\lceil\frac{\binom{h_1+k-2}{k-1}-1}{h_1-1}\right\rceil\right\}.
\end{equation}
The theorem provides an upper bound $\prod_jh_j$, not an assertion that
it matches this lower bound or solves diameter minimization.

\section{Algorithms, exact checks, and proof audit}
\label{sec:verification}

\subsection{Constructing and using a witness}

The witness is built by $r$ integer multiplications.  If
$L=\sum_j\lceil\log_2q_j\rceil$, every positive element has at most
$L+1$ binary digits, and the entire list has $O(rL)$ bits.  Even elementary
integer multiplication therefore gives a polynomial-time construction in
the bit length of the input list.  This is an existence algorithm; it does
not solve a shortest-vector problem for an arbitrary lattice.

Membership in $hA(q)$ requires $r$ Euclidean divisions to get the digits,
followed by checking $t+s(d)\le h$.  A canonical length-$h$ representation,
when it exists, is
\[
 (h-t-s(d),\ d_1,\ldots,d_r,\ t).
\]
Computing the whole digit-sum polynomial can be done using a sliding-window
convolution by $1+z+\cdots+z^{q_j-1}$ at each stage, requiring
$O(r(D+1))$ integer additions or subtractions.  Formula~\eqref{eq:digitscount}
then evaluates any $N(h)$ in $O(D+1)$ operations, or
\eqref{eq:affine} applies immediately beyond the exact threshold.
For small rank and very large radices,~\eqref{eq:inclusion} avoids creating
a degree-$D$ polynomial.

\subsection{What the verifier enumerates independently}

The main verifier does not infer minima from the claimed answer.  For a set
$A$ and a bound $H$, it enumerates every weak composition
$(m_0,\ldots,m_{k-1})$ of $H$ and groups them by their weighted sum.
For each group, it subtracts one representative from every other member.
These differences span exactly the degree-at-most-$H$ relation space.
Indeed, any relation of degree $d\le H$ can have its positive and negative
parts padded by $H-d$ zeros, so it appears as a difference in a group;
any difference of two group members is the difference of their two
representative differences.  The rank is computed by fraction-free integer
Gaussian elimination, not by floating-point singular values.

The recorded checks are summarized below.  Overlap between the ordered and
arbitrary-order profile collections is intentional; their counts are not
advertised as disjoint cases.
\begin{center}
\begin{tabular}{@{}p{0.69\linewidth}r@{}}
\toprule
\textbf{Check} & \textbf{Recorded count}\\
\midrule
Nondecreasing lists, $1\le r\le4$, $2\le q_j\le6$ & 125\\
Degree/rank checks for those lists & 754\\
Weak compositions enumerated for those lists & 56,493\\
Collision generators checked for those lists & 10,943\\
Ordered tuples in every order, $1\le r\le4$, $2\le q_j\le5$ & 340\\
Degree/rank checks for those tuples & 1,886\\
Weak compositions enumerated for those tuples & 109,392\\
Collision generators checked for those tuples & 22,909\\
Independent repeated-set-addition profiles & 34\\
Exact sumset-size checks by that method & 346\\
Individual membership tests against those sumsets & 134,080\\
Exhaustive vector-ball comparisons using all multiset pairs & 5\\
Vectors in those five balls, including their zeros & 79\\
\bottomrule
\end{tabular}
\end{center}
Every recorded check passed.  Additional examples include rank zero,
repeated minima, widely separated minima, and the two non-equivalent
sumset sequences in Section~\ref{sec:limits}.  All computations use Python
integers and the standard library; no external symbolic service or random
sampling of relations is involved.

In the five ball checks, every pair of representations in every bucket is
used, not only differences to a representative.  This independently gives
the entire finite relation ball and is compared with the bounded
carry-coordinate enumerator.  Entrance rigidity is checked on those balls.
The finite tests supplement the proof; they do not replace any quantifier
in Theorem~\ref{thm:spectrum}.

\subsection{The principal logical hazards}

A list of $r$ relations with the desired lengths proves only upper bounds
on the successive minima.  Lemma~\ref{lem:bound} supplies the missing lower
bounds for every integer combination.  Triangular elimination separately
establishes integral generation, rather than merely full rank.  Ties are
handled by the exact rank profile, with no assumption that the minima are
strictly increasing.

The normal-form proof counts integers only after uniqueness is established.
The ideal proof separately establishes kernel equality, which does not
follow from a lattice basis alone.  Finally, affine growth starts at $D-1$,
not merely at the sufficient but non-sharp bound $D$.

\section{Conclusion}

The realization problem has an explicit all-rank answer.  Prescribed
half-minima $h_j\ge2$ can be installed as successive multiplicative ratios
in a set containing zero and one.  The carry relations furnish the upper
bounds, and the partial weighted sums furnish simultaneous lower bounds
on every carry coordinate.  There are no additional arithmetic obstructions
on the list of minima.

The same construction supports more structure than the existence theorem
requires: its short-relation filtration is a direct-summand filtration,
its minimizing bases are rigid up to signs and ties, and its sumsets admit
unique digit normal forms.  The resulting product generating function and
sharp affine onset are consequences of those normal forms.  None of these
claims depends on an unproved general principle about arbitrary lattices
or arbitrary sumsets.

The mathematical proof and the bibliographic status are separate matters.
The former is fully given here.  The latter remains a research-draft claim
of a resolution of the conjecture as located in the inspected source;
independent scrutiny and a broader priority check remain appropriate.

\clearpage
\appendix
\section{The randomized area catalogue}\label{app:areas}

The catalogue was fixed before the draw.  The call was
\texttt{secrets.randbelow(96)}; its recorded zero-based output was $0$.
Thus the chosen area is entry~1.  The saved result, rather than a seeded
replay, is the audit record: operating-system entropy is not reproducible
from a seed that was never used.  The archived script defaults to displaying
the existing record and will not overwrite it with a fresh draw.

\begin{multicols}{3}
\footnotesize
\interlinepenalty=10000
\begin{enumerate}[leftmargin=2.3em,label=\arabic*.,itemsep=0.12em,parsep=0pt]
\item Additive combinatorics
\item Algebraic combinatorics
\item Analytic combinatorics
\item Extremal graph theory
\item Structural graph theory
\item Graph coloring
\item Spectral graph theory
\item Hypergraph theory
\item Ramsey theory
\item Design theory
\item Finite geometry
\item Matroid theory
\item Discrete geometry
\item Convex geometry
\item Geometry of numbers
\item Lattice polytopes and Ehrhart theory
\item Tiling theory
\item Combinatorial game theory
\item Permutation patterns
\item Combinatorics on words
\item Automata and formal languages
\item Symbolic dynamics
\item Substitution dynamical systems
\item Ergodic theory
\item Topological dynamics
\item Arithmetic dynamics
\item Complex dynamics
\item Real dynamical systems
\item Numerical semigroups
\item Commutative algebra
\item Monomial ideals
\item Combinatorial commutative algebra
\item Invariant theory
\item Finite group theory
\item Combinatorial group theory
\item Semigroup theory
\item Nonassociative algebra
\item Lie algebras
\item Representation theory
\item Category theory
\item Type theory
\item Universal algebra
\item Order theory
\item Lattice theory
\item Computability theory
\item Proof theory
\item Model theory
\item Set-theoretic topology
\item General topology
\item Knot theory
\item Low-dimensional topology
\item Algebraic topology
\item Persistent homology
\item Metric geometry
\item Graph limits
\item Probabilistic combinatorics
\item Discrete probability
\item Markov chains
\item Random walks
\item Information theory
\item Coding theory
\item Finite fields
\item Elementary number theory
\item Diophantine approximation
\item Diophantine equations
\item Multiplicative number theory
\item Partition theory
\item q-series
\item Special functions
\item Orthogonal polynomials
\item Approximation theory
\item Potential theory
\item Harmonic analysis
\item Fourier analysis on finite groups
\item Functional equations
\item Matrix theory
\item Nonnegative matrices
\item Tensor algebra
\item Polynomial inequalities
\item Moment problems
\item Optimization
\item Discrete optimization
\item Numerical analysis
\item Difference equations
\item Differential algebra
\item Differential equations
\item D-finite functions
\item Integer sequences
\item Enumerative topology
\item Geometric group theory
\item Computational geometry
\item Boolean functions
\item Algebraic statistics
\item Tropical geometry
\item Real algebraic geometry
\item Incidence geometry

\end{enumerate}
\end{multicols}

\clearpage
\section{Reproduction and package contents}

Run the finite checks from the package root with
\begin{lstlisting}
python3 code/verify.py
\end{lstlisting}
The program regenerates the CSV and JSON outputs under \texttt{data/}.
The two main mathematical routines can also be used directly:
\begin{lstlisting}[language=Python]
# Add this package's code/ directory to the Python import path:
import sys
sys.path.insert(0, "code")
from mixed_radix import construct, carry_basis, sumset_size
from mixed_radix import canonical_representation

q = (2, 3, 5)
assert construct(q) == (0, 1, 2, 6, 30)
assert sumset_size(q, 6) == 105
assert canonical_representation(q, 4, 35) == (0, 1, 2, 0, 1)
\end{lstlisting}
Python also has a standard-library module named \texttt{code}; avoid using
that name as a package import.  The supported import is
\texttt{from mixed\_radix import ...} after adding the directory.

Compile the article using \texttt{latexmk -pdf article.tex}, or run
\texttt{pdflatex article.tex} twice.  The package includes the complete
\LaTeX{} source, this PDF, the implementation, the verifier, the recorded
area choice, exact results, and separate source and proof audits.  No source
paper PDF, font files, or external dependencies are redistributed.

\begin{thebibliography}{9}
\bibitem{obryant}
Kevin O'Bryant,
\emph{On Nathanson's Triangular Number Phenomenon},
arXiv:2506.20836v2, 18 August 2025.
\url{https://arxiv.org/abs/2506.20836v2}.
The target is the unnumbered conjecture after Lemma~1,
Section~1.1, page~2.  Definitions and the adjacent statement were checked
in both the HTML and the displayed PDF page on 20 September 2026.

\bibitem{nathanson}
Melvyn B. Nathanson,
\emph{Explicit sumset sizes in additive number theory},
Canadian Mathematical Bulletin, published online 4 December 2025;
preprint arXiv:2505.05329.
\url{https://arxiv.org/abs/2505.05329}.
Cited for the surrounding range-of-sumset-sizes problem, not as an ingredient
in the proofs of this article.
\end{thebibliography}
\end{document}
